% Sample article for the Electronic Journal of Combinatorics.
%
% EJC papers *must* begin with the following two lines. 
\documentclass[12pt]{article}
\usepackage[amsmath]{e-jc}

% Omit the option "[amsmath]" if you don't want to use the amsmath package.

% GEOMETRY
% Please remove all commands that change parameters such as
%    margins or page sizes. The style file sets them.
% Don't use the fullpage package.

% PACKAGES
% Packages amssymb, amsthm and hyperref are already loaded. 
% We recommend the graphicx package for images:
\usepackage{graphicx}
% Forbidden:  fullpage, authblk, babel

% THEOREM ENVIRONMENTS
% Theorem-like environments that are declared in the style file are:
%   theorem, lemma, corollary, proposition, fact, observation, claim, note
%   definition, example, conjecture, open, problem, question, remark

% CLEVEREF PACKAGE
% The cleveref package is sensitive to the order that packages are loaded,
% and is incompatible with some. If you use \usepackage{cleveref} yourself,
% it will probably not work properly (lemmas will be shown as theorems, etc.).
% The best way is to add it as an option to package e-jc, for example:
%    \usepackage[amsmath,cleveref]{e-jc}
% To pass options to the cleveref package you can do it like this:
%    \def\cleverefoptions{capitalize,noabbrev}
%    \usepackage[amsmath,cleveref]{e-jc}
% Leave out amsmath if you aren't using it, but if you are using
% both amsmath and cleveref you have to do it this way.

% CHARACTER CODES
% Please do not use non-ascii characters in this file, but instead use the LaTex
% macros for characters with diacritical marks, such as G\"{o}del, R\'{e}nyi,
% Erd\H{o}s.  Don't use the package "babel".  Note that this is the opposite
% of the rule for the article metadata that you enter on the web page; sorry
% for the confusion!

% DATES
% Give the submission and acceptance dates in the format shown.
% The editors will insert the publication date in the third argument.
\dateline{Jan 1, 2024}{Jan 2, 2024}{TBD}

% SUBJECT CODES
% Give one or more subject codes separated by commas.
% Codes are available from http://www.ams.org/mathscinet/freeTools.html
\MSC{05C88, 05C89}

% COPYRIGHT AND LICENCE
% E-JC leaves copyright with authors.
% Uncomment exactly one of the following copyright statements.

%    One author:  ==========================
%\Copyright{The author.}
%\Copyright{The author. Released under the CC BY license (International 4.0).}
%\Copyright{The author. Released under the CC BY-ND license (International 4.0).}
%    More than one author: ===================
%\Copyright{The authors.}
%\Copyright{The authors. Released under the CC BY license (International 4.0).}
\Copyright{The authors. Released under the CC BY-ND license (International 4.0).}

% See https://creativecommons.org/licenses/ for a full explanation of the 
% Creative Commons licenses.
%
% We strongly recommend CC BY-ND or CC BY. Both these licenses allow others
% to freely distribute your work while giving credit to you. The difference
% between them is that CC BY-ND only allows distribution unchanged and in
% whole, while CC BY also allows remixing, tweaking and building upon your work.  

% TITLE
% If needed, include a line break (\\) at an appropriate place in the title.
% Footnotes on titles are not allowed.
\title{An elementary proof\\ of the reconstruction conjecture}

% AUTHORS
% Input authors, affiliations and addresses as follows.

% The following style is used if the authors don't all have the
% same affiliations.

\author{Jo Firstauthor\authornote{1}
\and
Jim Secondauthor\authornote{2}
\and
Janet Thirdauthor\authornote{1,2}
}

\authortext{1}{Department of Mathematics, University of Sutherland, Reay, Scotland
   (\email{j.firstauthor@uos.ac.uk}, \email{j.thirdauthor@uos.ac.uk}).}
\authortext{2}{Stochastics Department, Boyle County College,  Danville,
Kentucky, USA (\email{jsecond@bcc.edu}).}

% Use one \authortext for each affiliation, numbered 1,2,...  and for each
% author indicate which \authortexts apply.  The email address for each
% author goes in the first \authortext for that author.

% If all the authors have the same affiliations, don't use \authornote
% at all. Put the affiliations and email addresses into a single
% \authortext with empty first argument, like this:
%   \authortext{}{Affilation (\email{email1}, \email{email2}).}

% Note that other information such as grants go in the Acknowledgement
% section at the end of the paper, not here.  Don't use \thanks.

% E-JC does not customarily mark one author as "corresponding author"
% and we discourage this practice. However, if this is important for you,
% you may include this information in the \authortext section alongside
% the affiliation for that author.

\begin{document}

\maketitle

% ABSTRACT
% E-JC papers must include an abstract. The abstract should consist of a
% succinct statement of background followed by a listing of the principal
% new results that are to be found in the paper. The abstract should be
% informative, clear, and as complete as possible. Phrases like
% "we investigate..." or "we study..." should be kept to a minimum in
% favor of "we prove that..."  or "we show that...".  Do not include equation
% numbers, unexpanded citations (such as "[23]"), or any other references
% to things in the paper that are not defined in the abstract. The abstract
% may be distributed without the rest of the  paper so it must be entirely
% self-contained.  Try to include all words and phrases that someone
% might search for when looking for your paper.

\begin{abstract}
  The reconstruction conjecture states that the multiset of 
  vertex-deleted subgraphs of a graph determines the graph, provided
  it has at least 3 vertices.  This problem was independently introduced
  by Paul Kelly (1957) and Stanis\l aw Ulam (1960). In this paper,
  we prove the conjecture by elementary methods. 
  It is only necessary
  to integrate the Lenkle potential of the Broglington manifold over
  the quantum supervacillatory measure in order to reduce the set of
  possible counterexamples to a small number (less than a trillion).
  A simple computer program that implements Pipletti's classification
  theorem for torsion-free Aramaic groups with simplectic socles can
  then finish the remaining cases.
\end{abstract}

\section{Introduction}

The reconstruction conjecture states that the multiset of unlabeled
vertex-deleted subgraphs of a graph determines the graph, provided it
has at least three vertices.  This problem was independently introduced
by Kelly~\cite{Kelly} and Ulam~\cite{Ulam}.   The reconstruction
conjecture is widely studied
\cite{Bollobas,FGH,HHRT,KSU,Stockmeyer,WS} and is very interesting.
See \cite{BH} for more about the reconstruction conjecture.

\begin{definition} 
  A graph is \emph{fabulous} if \emph{rest of definition here}.
\end{definition}

\begin{theorem}\label{Thm:FabGraphs}
  All planar graphs are fabulous.
\end{theorem}

\begin{proof}
  Suppose on the contrary that some planar graph is not fabulous.
  Then by well-ordering there is a smallest planar graph that is
  not fabulous.  It is not the trivial graph, and we can easily see 
  that the property of being not fabulous is preserved by edge contraction.
  This gives a contradiction.
\end{proof}

%%%%%%%%%%%%%%%%%%%%%%%%%%%%%%%%%%%%%%%%%%%%%%%%%%%
\section{Broglington Manifolds}

This section describes background information about Broglington
Manifolds.

\begin{lemma}\label{lem:Technical}
  Broglington manifolds are abundant.
\end{lemma}

\begin{proof}
  A proof is given here.
\end{proof}

%%%%%%%%%%%%%%%%%%%%%%%%%%%%%%%%%%%%%%%%%%%%%%%%%
\section{Proof of Theorem~\ref{Thm:FabGraphs}}

In this section we complete the proof of Theorem~\ref{Thm:FabGraphs}.

\begin{proof}[Proof of Theorem~\ref{Thm:FabGraphs}]
Let $G$ be a graph. We have
 % The align environment for multi-line equations is defined in the amsmath package.
  \begin{align}
    |X| &= a+b+c \nonumber\\
         &= \alpha\beta\gamma.
  \end{align}
  This completes the proof of Theorem~\ref{Thm:FabGraphs}.
\end{proof}

\begin{figure}[ht]
  \centering
    % Use \includegraphics to import figures; for example 
    %      \includegraphics[scale=0.6]{filename}
  \caption{Here is an informative figure.\label{fig:InformativeFigure}}
\end{figure}

%%%%%%%%%%%%%%%%%%%%%%%%%%%%%%%%%%%%%%%%%%%%%%%%%
% The following optional unnumbered section is where you put personal acknowledgements,
% research grant support, and similar things.  Do not put them on the front page.
\subsection*{Acknowledgements}

Thanks to Professor Qwerty for suggesting the proof of
Lemma~\ref{lem:Technical}.
Jim Secondauthor received support from NASA grant MOON1969.

%BIBLIOGRAPHY
% You do not have to use the same format for your references, but 
%    include everything in this file.
% If you use BibTeX to create a bibliography, copy the .bbl file into here.
% We recommend you use \doi{...} and \arxiv{...} like the examples below,
% as they give a short display form with an active link to the full url.

\begin{thebibliography}{99}

\bibitem{Bollobas} B. Bollob{\'a}s.  Almost every
  graph has reconstruction number three.  \emph{J. Graph Theory},
  14(1): 1--4, 1990.

\bibitem{BH} J.\,A. Bondy and R. Hemminger,
 Graph reconstruction---a survey.
\emph{J. Graph Theory}, 1: 227--268, 1977. \doi{10.1002/jgt.3190010306}.

\bibitem{FGH} J.~Fisher, R.\,L. Graham, and F.~Harary.  
  A simpler counterexample to the reconstruction conjecture for
  denumerable graphs.  \emph{J. Combinatorial Theory, Ser. B},
  12: 203--204, 1972.

\bibitem{HHRT} E. Hemaspaandra, L.\,A. Hemaspaandra,
  S.~P. Radziszowski, and R. Tripathi. 
  Complexity results in graph reconstruction.  \emph{Discrete
    Appl. Math.}, 155(2): 103--118, 2007.

\bibitem{Kelly} P.\,J. Kelly.  A congruence theorem for
  trees.  \emph{Pacific J. Math.}, 7: 961--968, 1957.

\bibitem{KSU} M. Kiyomi, T. Saitoh, and R. Uehara.
   Reconstruction of interval graphs.  In 
    \emph{Computing and combinatorics}, volume 5609 of
    \emph{Lecture Notes in Comput. Sci.}, pages 106--115. Springer, 2009.

\bibitem{Stockmeyer} P.~K. Stockmeyer.  The falsity of the
  reconstruction conjecture for tournaments.  \emph{J. Graph
    Theory}, 1(1):19--25, 1977.

\bibitem{Ulam} S.\,M. Ulam.  A collection of
    mathematical problems.  Interscience Tracts in Pure and
  Applied Mathematics, no. 8.  Interscience Publishers, New
  York-London, 1960.
  
\bibitem{WS} D.~B. West and H. Spinoza.
  Reconstruction from $k$-decks for graphs with maximum degree~2.
  \arxiv{1609.00284vi}, 2016.

\end{thebibliography}

\end{document}
